\cleardoublepage
\section{Prime and Composite Numbers}

Some numbers stand alone. Others are made of parts. Prime numbers have only two factors: $1$ and themselves. Composite numbers have more. This is not just a label—it is a clue to how numbers work.

$2, 3, 5, 7, 11, 13,$ and $17$ are prime. They cannot be broken down. But $15$ is composite: it splits into $3$ and $5$. Learning to recognize prime and composite numbers helps us factor numbers, simplify fractions, and understand the building blocks of math.

\subsection{Example: Prime and Composite Numbers}

As you explore numbers further, you'll discover that some numbers have only two factors, while others have many. A \textbf{prime number} is a number greater than $1$ that has exactly two distinct factors: $1$ and itself. For example, $2, 3, 5, 7, 11, 13, 17,$ and $19$ are all prime numbers under $20$. In contrast, a \textbf{composite number} has more than two factors. For instance, $15$ is composite because it can be divided by $1, 3, 5,$ and $15$.

Let's practice identifying these. Consider the numbers $14, 17, 21, 23,$ and $25$. Among these, $17$ and $23$ are prime because they have only two factors. The others—$14, 21,$ and $25$—are composite. If you look at all the composite numbers between $10$ and $20$, you'll find $10, 12, 14, 15, 16, 18,$ and $20$. It's also important to note that the number $1$ is neither prime nor composite, since it has only one factor: itself. Recognizing prime and composite numbers is a foundational skill for understanding more advanced math concepts like prime factorization and divisibility.

\newpage 

\subsection{Prime Factorization Questions}

\begin{enumerate}
  \item \textbf{Prime Factorization (Multiple Choice)} \\ 
  What is the prime factorization of $60$?
  \begin{itemize}
    \item[A.] $2 \times 2 \times 3 \times 5$
    \item[B.] $2 \times 3 \times 10$
    \item[C.] $2 \times 5 \times 6$
    \item[D.] $3 \times 4 \times 5$
  \end{itemize}
  \textbf{Answer:} The correct answer is \textbf{A}. The prime factorization of $60$ is $2 \times 2 \times 3 \times 5$ (or $2^2 \times 3 \times 5$).

  \item \textbf{Prime Factorization (Open-Ended)} \\ 
  Find the prime factorization of $84$.
  
  \textbf{Answer:} The prime factorization of $84$ is $2 \times 2 \times 3 \times 7$ (or $2^2 \times 3 \times 7$).

  \item \textbf{Prime Factorization (Open-Ended)} \\ 
  Write the prime factorization of $105$ using exponents.
  
  \textbf{Answer:} The prime factorization of $105$ is $3 \times 5 \times 7$. Since all are single primes, no exponents are needed.

  \item \textbf{Prime Factorization (Multiple Choice)} \\ 
  Which of the following is the prime factorization of $72$?
  \begin{itemize}
    \item[A.] $2^3 \times 3^2$
    \item[B.] $2^2 \times 3^3$
    \item[C.] $2^4 \times 3$
    \item[D.] $2 \times 6^2$
  \end{itemize}
  \textbf{Answer:} The correct answer is \textbf{A}. $72 = 2^3 \times 3^2$.
\end{enumerate}

\newpage
\subsection{Hand-drawing 1}
This is where you will see a hand-drawn answer to the previous question in \textbf{English}.
\begin{figure}[htbp]
  \centering
  \includegraphics[width=\linewidth]{example-image} % TODO: Update caption for actual content
  \caption{A hand-drawn answer in English.}
\end{figure}

\newpage
\subsection{Hand-drawing 2}
This is where you will see a hand-drawn answer to the previous question in \textbf{Korean}.
\begin{figure}[htbp]
  \centering
  \includegraphics[width=\linewidth]{example-image} % TODO: Update caption for actual content
  \caption{A hand-drawn answer in Korean.}
\end{figure}
